\clearpage\begingroup\SourceLeaf{691b}
\ArithmeticHead{56}{L}{H. GEPPERT, ÜBER GAUSS’ ARBEITEN ZUR MECHANIK UND POTENTIALTHEORIE.}
\makeatletter\tagsleft@true\let\veqno\@@leqno\makeatother
\renewcommand{\thefootnote}{\arabic{footnote})}\setcounter{footnote}{113}
\begin{center}\textls[100]{2. Torsionsbestimmungen.}\end{center}
\noindent\hspace*{\parindent}In der \textit{Intensitas vis magneticae}\footnote{Werke V, S. 79—118, bes. 93—100. \textsc{Ostwalds} Klassiker, Band 53, 1894, bes. S. 19—27.} beschäftigt sich \textsc{Gauss} in den Art.
8—9 mit der Frage der Beeinflussung der Schwingungen einer Magnetnadel
durch die Torsion des Aufhängefadens. Sind \(D,E\) ein unterer bezw. oberer
Durchmesser dieses Fadens, die im ungedrillten Zustande einander parallel
sind und nach der magnetischen Beeinflussung der Nadel mit dem magnetischen
Meridian bezw. die Winkel \(u\) und \(v\) bilden, so ist das Torsionsmoment gleich
\(\Theta(v-u)\), worin \(\Theta\) den noch zu bestimmenden Torsionsmodul des Fadens bezeichnet.
Genauer gesagt, genügt es bei dieser Frage, den Quotienten
\[
n=\frac{HM}{\Theta}
\]
\noindent zu ermitteln, worin \(H\) die Horizontalintensität der erdmagnetischen Kraft, \(M\)
das magnetische Moment der Nadel bezeichnen. Die Werte für \(v\) können
beliebig eingestellt werden, die von \(u\) sind immer sehr klein. Man findet \(n\)
aus zwei Beobachtungen von korrespondierenden \(u,v\)-Werten nach der Formel
\[
n+1=\frac{v_2-v_1}{u_2-u_1}.\tag{4}
\]
\noindent Nun modifiziert aber die Torsion als Störungsglied auch die Schwingungsgleichung
der Nadel, und \textsc{Gauss} beweist aus deren Lösung, dass bei geringer
Abweichung der Ruhelage vom magnetischen Meridian die modifizierte Schwingungsdauer
\(T^*\) mit der idealen Schwingungsdauer \(T\) durch die Formel
\[
T=T^*\sqrt{\frac{n+1}{n}}
\]
\clearpage\SourceLeaf{692}\ArithmeticHead{57}{R}{V. ARBEITEN ZUR PRAKTISCHEN MECHANIK.}
\noindent verknüpft ist, mittels derer man den Einfluss der Torsion eliminieren kann.
Im Briefe an \textsc{Gerling} vom 15. Dez. 1835\footnote{\textsc{Gauss-Gerling}-Briefwechsel, S. 454—455.} gibt \textsc{Gauss} noch eine Anleitung
zur praktischen Bestimmung von \(n\). Man misst die Nadelstellung \(a\) auf
gewöhnliche Art, die Stellung \(b\), nachdem die Nadel um \(360^0\) rückwärts gedreht
ist, wieder die gewöhnliche Stellung \(a'\), die Stellung \(c\), nachdem die
Nadel um \(360^0\) vorwärts gedreht ist, und schliesslich wieder die gewöhnliche
Stellung \(a''\). Dann ist also
\[
n+1=\frac{2\pi}{b-a}=\frac{2\pi}{b-a'}=\frac{2\pi}{a'-c}=\frac{2\pi}{a''-c},
\]
\noindent und setzt man:
\[
m=\tfrac14(2b-2c-a+a''),
\]
\noindent so wird, da \(m\) gegen \(2\pi\) sehr klein ist,
\[
n\sim\frac{1296000''}{m''}.
\]
\noindent Diese Methode von \textsc{Gauss} ist sowohl von seinen Schülern\footnote{\textsc{Goldschmidt}, \textit{Über die Bestimmung der absoluten Intensität.} Resultate aus den Beobb. d. magn. Vereins 1840, S. 122—156, bes. 135—136.}, als auch von
\textsc{Lamont}\footnote{J. \textsc{Lamont}, \textit{Handbuch des Magnetismus}, Voss, Leipzig 1867. Erschienen in \textit{Allgemeine Encyklopädie der Physik}, herausgegeb. von \textsc{Gustav Karsten}. Bes. S. 132—133.} stets in Anwendung gebracht worden.

\begin{center}\textls[100]{3. Bestimmung von Trägheitsmomenten.}\end{center}
\noindent\hspace*{\parindent}Ist \(K\) das Trägheitsmoment der Magnetnadel um ihren Aufhängepunkt,
so gilt für ihre Schwingungsdauer die Relation
\[
T=\pi\sqrt{\frac{K}{HM}},
\]
\noindent und die Bestimmung von \(HM\) erfordert also die Kenntnis von \(K\). An der
gleichen Stelle, Art. 10—11, entwickelt \textsc{Gauss} ein Verfahren zur Ermittlung
von \(K\). Es besteht darin, dass mit der Nadel ein hölzerner Querstab verbunden
wird, an dem in zwei zur Aufhängung symmetrischen Punkten \(A,B\)
mit dem Abstande \(2r'\) zwei gleiche Gewichte \(p\) hängen; ist dann \(C\) das Trägheitsmoment
des Stabes bezüglich der Aufhängeachse, vermehrt um die Trägheitsmomente
der Gewichte \(p\) bezüglich der durch ihre Schwerpunkte und Aufhängelager
gehenden Vertikalen, so ist das gesamte zu \(K\) hinzutretende Trägheitsmoment
\par\noindent\textbf{X\textsubscript{2} Abh. 7.}\hfill 8
\clearpage\SourceLeaf{693}\ArithmeticHead{58}{L}{H. GEPPERT, ÜBER GAUSS’ ARBEITEN ZUR MECHANIK UND POTENTIALTHEORIE.}
\noindent\(C+2pr^{\prime2}\); ist \(T\) die Schwingungsdauer der unbelasteten, \(T_1\) die
der belasteten Nadel, so gilt
\[
\pi^2K=HMT^2;\quad\pi^2(K+C+2pr^{\prime2})=HMT_1^2.
\]
\noindent Verändert man dann \(r'\) in \(r''\), wobei \(T_1\) in \(T_2\) übergehen möge, so ist auch
\[
\pi^2(K+C+2pr^{\prime\prime2})=HMT_2^2,
\]
\noindent und man hat also drei Gleichungen, die man nach den drei Unbekannten
\(HM,K\) und \(C\) auflösen kann.

Diese \textsc{Gauss}sche Methode geht auf \textsc{Poisson} zurück. \textsc{Gauss} verwandte
bei seinen Versuchen Gewichte, die kleine Ösen mit feinen Stiften trugen,
die auf den Träger an vorausbezeichneten Stellen aufgesetzt wurden. Diese
\textsc{Gauss}sche Methode ist in ihren praktischen Anwendungen eingehender entwickelt
worden von \textsc{Goldschmidt}\footnote{Vergl. a.\,a.\,O. S. 131 ff.}, der aber schon bemerkte, dass durch die
drehende Bewegung der Gewichte die Nadelschwingungen gestört und ausserdem
die Trägheitsmomente der Gewichte vergrössert werden. An dieser Ungenauigkeit
leiden auch spätere Versuchsreihen von \textsc{Hansteen} und \textsc{Sartorius von Waltershausen}.
Man umging diese Schwierigkeit nach dem Vorgange von \textsc{W. Weber}
durch Anbringung von Schneiden statt Stiften, oder dadurch, dass man als
Gewichte durchbohrte Zylinder verwandte, die auf vertikale Schäfte aufgesetzt
wurden. Trotz dieser Kritik hat sich die \textsc{Gauss}sche Methode noch lange erhalten,
z.\,B. auch bei \textsc{Lamont}\footnote{a.\,a.\,O. 15) S. 15.}. \textsc{E. Dorn}\footnote{\textsc{E. Dorn}, \textit{Die Reduktion der Siemensschen Einheit auf absolutes Mass.} \textsc{Wiedemanns} Ann. d. Phys., 25, 1882, S. 773—816.} übernahm später die \textsc{Weber}sche
Methode und änderte sie in der Weise ab, dass nicht \(r'\), sondern \(p\) variiert
wird, d.\,h. er setzt auf den Querstab, der zwei symmetrisch zur Achse liegende
Zapfen trägt, verschiedene Gewichte auf, die die Form durchbohrter Zylinder
haben. \textsc{Kreichgauer}\footnote{\textsc{Kreichgauer}, \textit{Zur Bestimmung von Trägheitsmomenten durch Schwingungsversuch.} \textsc{Wiedemanns} Ann. d. Phys., 25, 1885, S. 273—308.} hat auf einen prinzipiellen Mangel der \textsc{Gauss}schen
und \textsc{Weber}schen Methode aufmerksam gemacht, der auf den nicht zu vermeidenden
elastischen Schwingungen des belasteten Systems beruht.
\clearpage\endgroup
